\centerline{\zihao{3}\bfseries ABSTRACT}

\vspace{1em}

The deepening of financial inclusion has driven rapid growth in the micro-loan industry, yet uneven service quality has emerged as a critical bottleneck constraining the sector\textquotesingle s healthy development. As a licensed micro-loan institution, P Company faces persistent service quality challenges, including low customer satisfaction, inefficient complaint handling, and suboptimal first-pass rates in business processes. The pressing question confronting P Company is how to systematically diagnose service quality deficiencies and implement effective improvements.

This study adopts the DMAIC methodology as its central framework and constructs a service quality improvement model tailored to the micro-loan context. In the Define phase, critical-to-quality characteristics were identified and improvement objectives established. In the Measure phase, the classic five-dimensional SERVQUAL model was extended to an eight-dimensional evaluation framework incorporating convenience, transparency, and product suitability. A Likert 5-point questionnaire was designed and administered, yielding 486 valid responses with strong reliability and validity (Cronbach\textquotesingle s $\alpha$ = 0.916, KMO = 0.894). In the Analyze phase, SERVQUAL gap analysis combined with Importance-Performance Analysis (IPA) revealed that the transparency dimension exhibited the largest perception-expectation gap ($-$1.70), with six key indicators falling within the IPA priority improvement zone. The Analytic Hierarchy Process (AHP) was then applied to structure the judgments of five industry experts (CR = 0.038, indicating satisfactory consistency), quantitatively determining improvement priorities. In the Improve phase, five targeted improvement initiatives were designed and implemented accordingly. In the Control phase, a key indicator monitoring mechanism was established to sustain the improvement outcomes.

Post-implementation results demonstrate substantial improvements across all service quality metrics: overall customer satisfaction rose from 3.12 to 3.85, the complaint closure rate increased from 62\% to 88\%, the first-pass rate improved from 62\% to 78\%, and average processing time was reduced from 12 working days to 7.5 working days. The findings confirm that the DMAIC methodology is applicable to micro-loan service quality improvement, and that the eight-dimensional evaluation model developed in this study possesses strong contextual suitability for micro-credit scenarios, offering a replicable improvement pathway for similar institutions.

\vspace{1em}
\noindent\textbf{\zihao{4} Keywords:} service quality; DMAIC; micro-loan; SERVQUAL; financial inclusion; IPA; AHP
